\section{Artefakte und Reproduktion}

\begin{table}[H]
  \centering
  \caption{Artefakte des Updates.}
  \label{tab:artefakte}
  \begin{tabular}{ll}
    \toprule
    Datei & Inhalt\\
    \midrule
    \texttt{UPDATE\_2026\_REPORT.md} / \texttt{.pdf} & dieser Report\\
    \texttt{figs/update\_2026\_a2c.png}  & End-NAV je Episode, Training vs.\ Test\\
    \texttt{figs/update\_2026\_1y.png}   & Kursentwicklung der drei Aktien, letzte zwölf Monate\\
    \texttt{updated\_metrics\_final.json} & aktualisierte Fundamentalkennzahlen\\
    \texttt{results\_ba\_newdata.json}   & Roh-Ergebnisse der fünf Agenten (Training + Test)\\
    \texttt{agent\_ba\_1..5.zip}         & die fünf neu trainierten A2C-Modelle\\
    \texttt{env\_runner.py}              & Environment-Aufsatz (lokale Daten, Seriell-Patch)\\
    \texttt{train\_eval.py}              & Training und Auswertung (Notebook-Protokoll)\\
    \texttt{plot\_results.py}            & Tabellen-Ausgabe und Abbildung\\
    \texttt{BA\_5Y\_daily.csv}           & Boeing-Tagesdaten (fünf Jahre)\\
    \bottomrule
  \end{tabular}
\end{table}

\paragraph{Reproduzieren.}
{\footnotesize
\begin{verbatim}
cd aif-latex
python3 aif_run/train_eval.py      # trainiert 5 Agenten + wertet aus (~35 min)
python3 aif_run/plot_results.py    # Tabelle + Abbildung
python3 aif_run/make_pdf.py UPDATE_2026_REPORT.md AIF_Update_2026.pdf
\end{verbatim}
}

\paragraph{Hinweis zu den Datenquellen.} Yahoo/\texttt{yfinance} ist von diesem Host hart
geblockt (HTTP 429 auf allen Endpunkten, auch nach zwanzig Minuten). Boeing und Toyota
stammen daher aus der History-API von stockanalysis.com, die Airbus-Kurse aus ariva.de
sowie der EPA:AIR-Seite von stockanalysis.com. Achtung: der Endpunkt
\texttt{/api/symbol/s/AIR/history} liefert \emph{AAR Corp (US)} und nicht Airbus -- die
daraus zunächst abgeleiteten Airbus-Zahlen waren falsch und wurden verworfen.
